%
% voice-call.tex
%
% Z specification of the voice call in voxd
% (src/punt_vox/voxd/voice_call/, PL-1371): the call lifecycle that
% call_start, call_stop and call_status drive, the playback ownership of
% AgentSpeaker, and the per-call tool result cache. The tests in
% tests/voice_call/ assert the named properties.
%
% PL-1779 adds the session handoff: a briefed call collects the person's
% answers and decisions while it is live and seals them as complete or
% partial before it settles. The design is docs/voice-call-handoff.md.
%

\documentclass[a4paper,10pt,fleqn]{article}
\usepackage[margin=1in]{geometry}
\usepackage{fuzz}

\begin{document}

\title{Vox Call: A Z Specification}
\author{Punt Labs, for PL-1371 and PL-1779}
\date{November 2026}
\maketitle

\section{Introduction}

voxd runs at most one call. \texttt{call\_start} creates it in
\emph{connecting}; the call task provisions the ElevenLabs agent and opens
the conversation, which makes it \emph{live}. A remote hang-up, the agent's
built-in \texttt{end\_call} tool or a \texttt{call\_stop} ends it; a connect timeout or other error fails it.
While a call holds the output, it owns the playback mutex across replies;
an interruption ends the current reply only. Each tool call is keyed by its
canonical arguments; a durable key runs its operation at most once per call.

A call started with a brief from a coding session also collects a
handoff: the person's answers, decisions and instructions, noted while the
call is live and sealed when it ends. The handoff section names nine more
properties, H1 to H9, and leaves C1 to C12 as they are.

The model names twelve call properties:

\begin{enumerate}
  \item[C1] \emph{One call}: a second \texttt{call\_start} is refused while a
    call is connecting or live (\texttt{test\_call\_barge\_in\_reask\_starts\_one\_worker}
    asserts it through the handlers).
  \item[C2] \emph{Stoppable at once}: \texttt{call\_stop} ends a call that is
    still provisioning or connecting
    (\texttt{test\_stop\_cancels\_a\_call\_still\_provisioning},
    \texttt{test\_stop\_cancels\_a\_call\_still\_connecting}).
  \item[C3] \emph{Terminal states stay}: an ended or failed call never
    becomes live again.
  \item[C4] \emph{Output held across a cut}: an interruption ends the reply
    but not the call's hold on playback
    (\texttt{test\_the\_call\_keeps\_the\_output\_across\_a\_cut}).
  \item[C5] \emph{Output released before cleanup finishes}: a remote hang-up
    can drain the last reply, but no settled call holds playback.
  \item[C6] \emph{Durable once}: a re-asked durable tool call replays and
    never runs its operation a second time
    (\texttt{test\_reask\_during\_a\_slow\_start\_starts\_one\_worker}).
  \item[C7] \emph{Return opens with only away results and open decisions}:
    compare the registry snapshot with the previous live call's durable departure
    cursor; a failed dial does not advance it.
  \item[C8] \emph{Return snapshot covers pre-cut events without losing later ones}:
    subscribe before taking the snapshot and suppress live reports already
    covered by it.
  \item[C9] \emph{Managed agent teardown is idle-only}: teardown does not
    end a call or alter its state
    (\texttt{test\_agent\_teardown\_refuses\_a\_live\_call}).
  \item[C10] \emph{Echo protection is declared before dialing}: microphone
    calls without a headset or external AEC declaration are refused before
    provisioning; scripted calls do not open the microphone
    (\texttt{test\_mic\_call\_without\_echo\_protection\_does\_not\_dial}).
  \item[C11] \emph{A named stop affects only that call}: a stop naming a
    previous call leaves its replacement unchanged
    (\texttt{test\_stale\_stop\_does\_not\_hang\_up\_the\_replacement\_call}).
  \item[C12] \emph{Attribution survives the close}: the agent's end\_call
    acknowledgement, taken while the call is live, names the agent as who
    ended it whatever end follows; a failure names nobody
    (\texttt{test\_the\_agent\_ending\_the\_call\_is\_credited\_to\_the\_agent},
    \texttt{test\_a\_stop\_after\_the\_agents\_end\_call\_is\_still\_the\_agents\_hang\_up},
    \texttt{test\_a\_failure\_after\_the\_agents\_end\_call\_is\_nobodys\_hang\_up}).
\end{enumerate}

\section{Types}

\begin{zed}
  [KEY, CALLID] \\
  CALL ::= idle | connecting | live | ended | failed \\
  INPUT ::= scripted | headset | externalAEC | openSpeakers \\
  ENDEDBY ::= nobody | byAgent | byPeer | byCaller \\
  ZBOOL ::= ztrue | zfalse
\end{zed}

\texttt{idle} is the handlers' state before the first call. \texttt{ENDEDBY}
is who ended the call: the agent through its built-in end\_call tool, the
peer closing the conversation, or the caller stopping it.

\section{State}

\begin{schema}{Call}
  call : CALL \\
  holding, replying, settled : ZBOOL \\
  durable, ran : \power KEY \\
  current, issued : \power CALLID \\
  endedBy : ENDEDBY
\where
  (call \in \{ idle, connecting \} \lor settled = ztrue) \implies holding = zfalse \\
  call = idle \implies settled = ztrue \\
  call \in \{ connecting, live \} \implies settled = zfalse \\
  replying = ztrue \implies holding = ztrue \\
  ran \subseteq durable \\
  current \subseteq issued \\
  call = idle \implies current = \emptyset \\
  call \neq idle \implies \# current = 1 \\
  call \in \{ idle, connecting, failed \} \implies endedBy = nobody \\
  call = live \implies endedBy \in \{ nobody, byAgent \} \\
  call = ended \implies endedBy \neq nobody
\end{schema}

Here \texttt{current} retains the last call's ID even after it ends, and
\texttt{issued} models the IDs reserved by the call trace store.

The first predicate is C5. \texttt{durable} is the set of durable keys whose
operation has run or is in flight; \texttt{ran} those that finished. The
last three predicates are C12: \texttt{endedBy} is set before the terminal
state only by the agent's end\_call acknowledgement, every ended call names
who ended it, and a failed call names nobody.

\begin{schema}{InitCall}
  Call' \\
  savedIds? : \power CALLID
\where
  call' = idle \\
  holding' = zfalse \\
  replying' = zfalse \\
  settled' = ztrue \\
  durable' = \emptyset \\
  ran' = \emptyset \\
  current' = \emptyset \\
  issued' = savedIds? \\
  endedBy' = nobody
\end{schema}

\section{Lifecycle}

\begin{schema}{ManagedAgent}
  registered, provisioning, awaitingProvision : ZBOOL
\where
  registered = ztrue \implies provisioning = zfalse \\
  awaitingProvision = ztrue \implies provisioning = ztrue
\end{schema}

\begin{schema}{InitManagedAgent}
  ManagedAgent' \\
  savedRegistration? : ZBOOL
\where
  registered' = savedRegistration? \\
  provisioning' = zfalse \\
  awaitingProvision' = zfalse
\end{schema}

\texttt{registered} denotes a complete on-disk registration. It may
predate the daemon, which does not discard it. Partially created resources
are represented only as retryable cleanup: the next provisioning attempt
cleans the partial record first, under the same registration lock.

\begin{schema}{Start}
  \Delta Call \\
  \Xi ManagedAgent \\
  input? : INPUT \\
  id? : CALLID
\where
  call \in \{ idle, ended, failed \} \\
  settled = ztrue \\
  input? \neq openSpeakers \\
  id? \notin issued \\
  call' = connecting \\
  settled' = zfalse \\
  holding' = zfalse \\
  replying' = zfalse \\
  durable' = \emptyset \\
  ran' = \emptyset \\
  current' = \{ id? \} \\
  issued' = issued \cup current' \\
  endedBy' = nobody
\end{schema}

\texttt{Start} is refused until the prior call task has finished (C1).
A new call gets a fresh cache. A request that fails before launching a
task leaves the previous call unchanged.

\begin{schema}{BeginProvisionAgent}
  \Xi Call \\
  \Delta ManagedAgent
\where
  call = connecting \\
  registered = zfalse \\
  provisioning = zfalse \\
  registered' = registered \\
  provisioning' = ztrue \\
  awaitingProvision' = ztrue
\end{schema}

\begin{schema}{ProvisionAgent}
  \Delta Call \\
  \Delta ManagedAgent \\
  success? : ZBOOL
\where
  provisioning = ztrue \\
  registered' = success? \\
  provisioning' = zfalse \\
  awaitingProvision' = zfalse \\
  (awaitingProvision = ztrue \land success? = zfalse) \implies
    (call = connecting \land call' = failed) \\
  (awaitingProvision = zfalse \lor success? = ztrue) \implies call' = call \\
  settled' = settled \\
  holding' = holding \\
  replying' = replying \\
  durable' = durable \\
  ran' = ran \\
  current' = current \\
  issued' = issued \\
  endedBy' = endedBy
\end{schema}

Cancellation stops the call's wait for the provisioning thread, not the
thread itself. \texttt{ProvisionAgent} records whether registration succeeded,
not whether the call connected. Failure of the attempt awaited by the
connecting call fails that call. Once detached, the attempt can finish
without changing any call state, including that of a later call.

\begin{schema}{Connect}
  \Delta Call \\
  \Xi ManagedAgent \\
  accepted? : ZBOOL
\where
  call = connecting \\
  registered = ztrue \\
  accepted? = ztrue \\
  call' = live \\
  settled' = settled \\
  holding' = holding \\
  replying' = replying \\
  durable' = durable \\
  ran' = ran \\
  current' = current \\
  issued' = issued \\
  endedBy' = endedBy
\end{schema}

\begin{schema}{AgentEndCall}
  \Delta Call \\
  \Xi ManagedAgent
\where
  call = live \\
  endedBy = nobody \\
  endedBy' = byAgent \\
  call' = call \\
  settled' = settled \\
  holding' = holding \\
  replying' = replying \\
  durable' = durable \\
  ran' = ran \\
  current' = current \\
  issued' = issued
\end{schema}

\texttt{AgentEndCall} is the agent's end\_call acknowledgement: the call
stays live, playing out the closing sentence, and already names the agent
as who ended it. The provider's close follows, through \texttt{End}.

Acceptance here includes validating or provisioning the registration,
obtaining the signed URL, preparing the worker snapshot, and completing the
ElevenLabs handshake. A mismatched registration is refused rather than
deleted on connect; failures outside registration itself end through
\texttt{End}. A connect timeout takes \texttt{End} with
$failure? = ztrue$ and $draining? = zfalse$. It cancels the task and
reports the failed status and error in both the start reply and later
status replies.

\begin{schema}{End}
  \Delta Call \\
  \Delta ManagedAgent \\
  failure?, draining? : ZBOOL \\
  by? : ENDEDBY
\where
  call \in \{ connecting, live \} \\
  by? \in \{ byPeer, byCaller \} \\
  failure? = ztrue \implies call' = failed \land endedBy' = nobody \\
  failure? = zfalse \implies call' = ended \\
  (failure? = zfalse \land endedBy = byAgent) \implies endedBy' = byAgent \\
  (failure? = zfalse \land endedBy = nobody) \implies endedBy' = by? \\
  settled' = zfalse \\
  draining? = ztrue \implies
    (call = live \land failure? = zfalse \land holding' = holding) \\
  draining? = zfalse \implies holding' = zfalse \\
  replying' = zfalse \\
  durable' = durable \\
  ran' = ran \\
  current' = current \\
  issued' = issued \\
  registered' = registered \\
  provisioning' = provisioning \\
  awaitingProvision' = zfalse
\end{schema}

\texttt{End} covers hang-up, stop, timeout and failure. Its precondition
admits \emph{connecting}, which is C2; no operation leaves \emph{ended} or
\emph{failed} except a new \texttt{Start}, which is C3. Terminal status
can be visible before the call task has finished its cleanup, including
draining the final reply after a remote hang-up. \texttt{by?} is the peer
for a remote close and the caller for a stop or cancellation; it names who
ended the call only when the agent had not already (C12). A failure, even
one after the agent's end\_call, names nobody.

\begin{schema}{EndNamed}
  End \\
  id? : CALLID
\where
  id? \in current \\
  failure? = zfalse \\
  draining? = zfalse \\
  by? = byCaller
\end{schema}

\begin{schema}{RejectStaleStop}
  \Xi Call \\
  \Xi ManagedAgent \\
  id? : CALLID
\where
  current \neq \emptyset \\
  id? \notin current
\end{schema}

\begin{schema}{StopNamedNoChange}
  \Xi Call \\
  \Xi ManagedAgent \\
  id? : CALLID
\where
  settled = ztrue \\
  current = \emptyset \lor id? \in current
\end{schema}

\begin{schema}{SettleCall}
  \Delta Call \\
  \Xi ManagedAgent
\where
  call \in \{ ended, failed \} \\
  settled = zfalse \\
  settled' = ztrue \\
  call' = call \\
  holding' = zfalse \\
  replying' = zfalse \\
  durable' = durable \\
  ran' = ran \\
  current' = current \\
  issued' = issued \\
  endedBy' = endedBy
\end{schema}

\begin{zed}
  StopNamed \defs (EndNamed \hide (failure?, draining?, by?)) \semi SettleCall
\end{zed}

\begin{schema}{StopNamedDuringCleanup}
  SettleCall \\
  id? : CALLID
\where
  id? \in current
\end{schema}

A named stop checks identity before cancellation, under the same lock as
\texttt{Start}. A mismatch leaves the call unchanged (C11).
\texttt{StopNamed} ends a connecting or live call and settles it before
the stop replies. If the call has already ended or failed but its task is
still cleaning up, a matching stop cancels the remaining work and waits
for settlement without changing the terminal status. Before the first
call or after the matching call has settled, stopping is a no-op. A stop
without an ID targets the current call.
The stop schemas describe the state when the reply is sent, not the
internal task cancellation and wait.

\begin{schema}{TeardownAgent}
  \Xi Call \\
  \Delta ManagedAgent
\where
  call \in \{ idle, ended, failed \} \\
  settled = ztrue \\
  provisioning = zfalse \\
  registered' = zfalse \\
  provisioning' = provisioning \\
  awaitingProvision' = awaitingProvision
\end{schema}

\begin{schema}{TeardownRefused}
  \Xi Call \\
  \Xi ManagedAgent
\where
  settled = zfalse \lor provisioning = ztrue
\end{schema}

\begin{schema}{TeardownIncomplete}
  \Xi Call \\
  \Delta ManagedAgent \\
  stillRegistered? : ZBOOL
\where
  settled = ztrue \\
  provisioning = zfalse \\
  stillRegistered? = ztrue \implies registered = ztrue \\
  registered' = stillRegistered? \\
  provisioning' = provisioning \\
  awaitingProvision' = awaitingProvision
\end{schema}

\texttt{TeardownAgent} models success, including when nothing is registered.
The handlers serialize teardown with call startup and refuse it until the
call task finishes (C9). The registration lock then serializes the provider
changes: an in-flight provisioning attempt must finish before deletion can
succeed, or lock acquisition times out and \texttt{TeardownRefused} leaves
state unchanged. The model does not track the wait queue. A deletion
failure reports \texttt{TeardownIncomplete}; the surviving record may still
name the agent or may contain only tools to retry. No teardown outcome
touches workers or the departure cursor.

\section{Playback}

\begin{schema}{Speak}
  \Delta Call
\where
  call = live \\
  holding' = ztrue \\
  replying' = ztrue \\
  call' = call \\
  settled' = settled \\
  durable' = durable \\
  ran' = ran \\
  current' = current \\
  issued' = issued \\
  endedBy' = endedBy
\end{schema}

\begin{schema}{Interrupt}
  \Delta Call
\where
  call = live \\
  replying' = zfalse \\
  holding' = holding \\
  call' = call \\
  settled' = settled \\
  durable' = durable \\
  ran' = ran \\
  current' = current \\
  issued' = issued \\
  endedBy' = endedBy
\end{schema}

\texttt{Interrupt} leaves \texttt{holding} as it was: C4.

\section{Tool results}

\begin{schema}{RunDurable}
  \Delta Call \\
  k? : KEY
\where
  call = live \\
  k? \notin durable \\
  durable' = durable \cup \{ k? \} \\
  ran' = ran \\
  call' = call \\
  settled' = settled \\
  holding' = holding \\
  replying' = replying \\
  current' = current \\
  issued' = issued \\
  endedBy' = endedBy
\end{schema}

\begin{schema}{Replay}
  \Xi Call \\
  k? : KEY
\where
  call = live \\
  k? \in durable
\end{schema}

\begin{schema}{Finish}
  \Delta Call \\
  k? : KEY
\where
  k? \in durable \setminus ran \\
  ran' = ran \cup \{ k? \} \\
  durable' = durable \\
  call' = call \\
  settled' = settled \\
  holding' = holding \\
  replying' = replying \\
  current' = current \\
  issued' = issued \\
  endedBy' = endedBy
\end{schema}

A key enters \texttt{durable} only through \texttt{RunDurable}, whose
precondition is $k? \notin durable$; a re-ask finds it there and replays,
whether the first run finished or is still in flight under the shield: C6.

\section{Return and live updates}

The durable departure cursor contains only result identities. Questions
come from the current projection, never from that cursor. The snapshot
and cut come from the same worker view, after subscription and before
the provider can speak. The provider gets one first-message override
only when the selected reports are nonempty.
Without a valid departure cursor, the call seeds the result cursor at dial
without an opener and saves the new cursor on departure. Subsequent
calls advance it for results in the opener or successfully submitted for
live speech, not for hidden context or fresh registry writes at hang-up.
Submission is not a receipt that the caller heard the playback.
Local hang-up stops playback; remote hang-up drains remaining audio.
Neither updates the cursor from a late snapshot. The live stream retains
its bounded queue; the cut filters duplicates among events that remain queued.

\begin{zed}
  [CALLWORKER]
\end{zed}

\begin{schema}{ReturnBoundary}
  departed : ZBOOL \\
  cursor, cut, known : CALLWORKER \pfun \nat \\
  subscribed, snapshotted : ZBOOL
\where
  departed = zfalse \implies cursor = \emptyset \\
  snapshotted = ztrue \implies subscribed = ztrue \\
  snapshotted = zfalse \implies cut = \emptyset \land known = \emptyset
\end{schema}

\begin{schema}{SubscribeForReturn}
  \Delta ReturnBoundary
\where
  subscribed = zfalse \\
  subscribed' = ztrue \\
  snapshotted' = zfalse \\
  cut' = \emptyset \\
  known' = \emptyset \\
  departed' = departed \\
  cursor' = cursor
\end{schema}

\begin{schema}{ReturnSnapshot}
  \Delta ReturnBoundary \\
  newest?, result? : CALLWORKER \pfun \nat \\
  questions? : \power CALLWORKER \\
  newResults!, openQuestions! : \power CALLWORKER
\where
  subscribed = ztrue \\
  snapshotted = zfalse \\
  snapshotted' = ztrue \\
  subscribed' = subscribed \\
  departed' = departed \\
  cursor' = cursor \\
  cut' = newest? \\
  departed = zfalse \implies
    known' = result? \land
    newResults! = \emptyset \land openQuestions! = \emptyset \\
  departed = ztrue \implies openQuestions! = questions? \\
  departed = ztrue \implies newResults! =
    \{ w : \dom result? | w \notin \dom cursor \lor result?~w > cursor~w \} \\
  departed = ztrue \implies known' = cursor \oplus (newResults! \dres result?)
\end{schema}

\begin{schema}{LiveUpdateAfterCut}
  \Xi ReturnBoundary \\
  worker? : CALLWORKER \\
  sequence? : \nat
\where
  subscribed = ztrue \\
  snapshotted = ztrue \\
  worker? \notin \dom cut \lor sequence? > cut~worker?
\end{schema}

\begin{schema}{SubmitLiveResult}
  \Delta ReturnBoundary \\
  worker? : CALLWORKER \\
  sequence? : \nat \\
  nonempty? : ZBOOL
\where
  subscribed = ztrue \\
  snapshotted = ztrue \\
  nonempty? = ztrue \\
  worker? \notin \dom cut \lor sequence? > cut~worker? \\
  (worker? \in \dom known \land sequence? \leq known~worker?) \implies known' = known \\
  (worker? \notin \dom known \lor sequence? > known~worker?) \implies
    known' = known \oplus \{ worker? \mapsto sequence? \} \\
  subscribed' = subscribed \\
  snapshotted' = snapshotted \\
  departed' = departed \\
  cursor' = cursor \\
  cut' = cut
\end{schema}

\begin{schema}{DepartCall}
  \Delta ReturnBoundary \\
  established? : ZBOOL
\where
  subscribed' = zfalse \\
  snapshotted' = zfalse \\
  cut' = \emptyset \\
  known' = \emptyset \\
  established? = ztrue \implies
    snapshotted = ztrue \land departed' = ztrue \land cursor' = known \\
  established? = zfalse \implies departed' = departed \land cursor' = cursor
\end{schema}

\section{Session handoff}

A call started from a coding session may carry a brief: a distilled
description of the session's situation and the questions it needs
answered. The agent asks them, notes each answer, decision and
instruction while the call is live, and the daemon seals the notes as a
\emph{handoff} when the call ends. The handoff is \emph{complete} only when
the agent hung up after every question was resolved; every other ending of
a briefed call, including a failure, seals it \emph{partial}, so the
session is told that the record may be short.

\begin{zed}
  [QUESTION, NOTE] \\
  HANDOFF ::= noBrief | collecting | complete | partial
\end{zed}

\texttt{noBrief} is the call as it was before the handoff: no brief, no
questions, no notes. \texttt{collecting} is a briefed call that has not
sealed. \texttt{QUESTION} is a question the brief asked; \texttt{NOTE} is
one recorded item (an answer, decision, instruction or open point).

\begin{schema}{HandoffRecord}
  handoff : HANDOFF \\
  tainted : ZBOOL \\
  asked, resolved : \power QUESTION \\
  noted, pending, held : \power NOTE
\where
  resolved \subseteq asked \\
  pending \cap noted = \emptyset \\
  held \subseteq pending \\
  handoff = noBrief \implies
    (asked = \emptyset \land noted = \emptyset \land pending = \emptyset
      \land held = \emptyset) \\
  handoff \neq noBrief \implies tainted = ztrue
\end{schema}

\texttt{resolved} holds the questions the person answered or deferred
aloud; a question never discussed stays in \texttt{asked} alone and the
daemon lists it as unresolved when it seals. \texttt{pending} holds notes
the agent has sent that are not recorded yet: in flight, or held for the
person's spoken yes. \texttt{tainted} is the call's provenance flag
(untrusted text has reached the conversation). A briefed call is tainted
from its start, because the brief is written by a model that may have read
untrusted text; the last predicate says so. The flag is modeled only for
briefed calls.

\begin{schema}{Briefed}
  Call \\
  HandoffRecord
\where
  call = idle \implies handoff = noBrief \\
  call \in \{ connecting, live \} \implies handoff \in \{ noBrief, collecting \} \\
  handoff \in \{ complete, partial \} \implies call \in \{ ended, failed \} \\
  settled = ztrue \implies handoff \neq collecting \\
  handoff = complete \implies
    (call = ended \land endedBy = byAgent \land asked \subseteq resolved
      \land pending = \emptyset) \\
  call = failed \implies handoff \neq complete
\end{schema}

The last three predicates are H3 to H5; a failed call that has not yet
sealed is still \emph{collecting}, and the seal turns it partial. A terminal call can show
\emph{collecting} between \texttt{End} and the seal, as it can show
\texttt{settled = zfalse} between \texttt{End} and \texttt{SettleCall};
the fourth predicate is why the seal comes first.

The handoff adds nine properties:

\begin{enumerate}
  \item[H1] \emph{Starts as asked}: a call started with a brief begins
    \emph{collecting} with every question asked and nothing resolved or
    noted; a call started without one begins \emph{noBrief}, the call as
    before.
  \item[H2] \emph{Notes only while live}: a note is received and recorded
    only on a live, collecting call, and only against questions the brief
    asked.
  \item[H3] \emph{Sealed before settled}: no settled call is still
    collecting, so a status reply that shows a settled call shows a sealed
    handoff.
  \item[H4] \emph{Complete is earned}: a complete handoff means the agent
    ended the call (C12), every question was resolved and no note is still
    pending. A brief with no questions is allowed; its handoff is complete
    when the agent ends the call, and may hold no notes.
  \item[H5] \emph{Partial is stated}: a briefed call that ended without
    earning complete, or failed, is partial; a failed call is never
    complete.
  \item[H6] \emph{Notes survive the cut}: sealing keeps every note made
    while the call was live, and every note still pending, whatever ended
    the call.
  \item[H7] \emph{Durable once}: a re-asked note replays and adds no second
    note (C6).
  \item[H8] \emph{A sealed handoff stays}: nothing but the next
    \texttt{Start} changes a complete or partial handoff (C3).
  \item[H9] \emph{Acting notes need a yes}: on a tainted call, a decision or
    an instruction is recorded only after the person's spoken yes; every
    briefed call is tainted, so this holds for the whole call.
\end{enumerate}

\begin{schema}{InitHandoff}
  HandoffRecord'
\where
  handoff' = noBrief \\
  tainted' = zfalse \\
  asked' = \emptyset \\
  resolved' = \emptyset \\
  noted' = \emptyset \\
  pending' = \emptyset \\
  held' = \emptyset
\end{schema}

\begin{schema}{BeginHandoff}
  \Delta HandoffRecord \\
  questions? : \power QUESTION
\where
  handoff' = collecting \\
  tainted' = ztrue \\
  asked' = questions? \\
  resolved' = \emptyset \\
  noted' = \emptyset \\
  pending' = \emptyset \\
  held' = \emptyset
\end{schema}

\begin{schema}{BeginNoHandoff}
  \Delta HandoffRecord
\where
  handoff' = noBrief \\
  tainted' = zfalse \\
  asked' = \emptyset \\
  resolved' = \emptyset \\
  noted' = \emptyset \\
  pending' = \emptyset \\
  held' = \emptyset
\end{schema}

\begin{schema}{ReceiveNote}
  \Delta HandoffRecord \\
  Call \\
  note? : NOTE
\where
  call = live \\
  handoff = collecting \\
  note? \notin noted \cup pending \\
  pending' = pending \cup \{ note? \} \\
  held' = held \\
  noted' = noted \\
  asked' = asked \\
  resolved' = resolved \\
  handoff' = handoff \\
  tainted' = tainted
\end{schema}

\texttt{ReceiveNote} is the daemon seeing the agent's \texttt{note\_handoff}
call, before it is validated, gated or recorded. It happens in the call's
listener as the tool call arrives, so a note the agent sent is never lost to
the cancellation that follows a hang-up: if the call ends before the note is
recorded, it is still in \texttt{pending} when the handoff seals.

\begin{schema}{NoteHandoff}
  \Delta HandoffRecord \\
  Call \\
  note? : NOTE \\
  answers? : \power QUESTION \\
  acting?, confirmed? : ZBOOL
\where
  call = live \\
  handoff = collecting \\
  note? \in pending \\
  answers? \subseteq asked \\
  (acting? = ztrue \land tainted = ztrue) \implies
    (confirmed? = ztrue \land note? \in held) \\
  noted' = noted \cup \{ note? \} \\
  pending' = pending \setminus \{ note? \} \\
  held' = held \setminus \{ note? \} \\
  resolved' = resolved \cup answers? \\
  asked' = asked \\
  handoff' = handoff \\
  tainted' = tainted
\end{schema}

\texttt{answers?} is the questions this note resolves: one for an answer
or a deferral, none for a decision, an instruction or an open point.
\texttt{acting?} marks an answer, a decision or an instruction, that is,
every note the session will act on; \texttt{confirmed?} is the person's
spoken yes to the daemon's read-back of it (H9). Only an open point is not
acting. The daemon also checks an answer against the question and its offered
options.

\begin{schema}{HoldNote}
  \Delta HandoffRecord \\
  Call \\
  note? : NOTE \\
  acting? : ZBOOL
\where
  call = live \\
  handoff = collecting \\
  note? \in pending \setminus held \\
  acting? = ztrue \land tainted = ztrue \\
  held' = held \cup \{ note? \} \\
  pending' = pending \\
  noted' = noted \\
  asked' = asked \\
  resolved' = resolved \\
  handoff' = handoff \\
  tainted' = tainted
\end{schema}

\texttt{HoldNote} is the confirmation gate holding an acting note on a
tainted call for the person's spoken yes: the note stays pending and is now
also \texttt{held}. A note the gate has not seen, and an open point, is never
held.

\begin{schema}{DeclineNote}
  \Delta HandoffRecord \\
  Call \\
  note? : NOTE
\where
  call = live \\
  handoff = collecting \\
  note? \in held \\
  held' = held \setminus \{ note? \} \\
  pending' = pending \setminus \{ note? \} \\
  noted' = noted \\
  asked' = asked \\
  resolved' = resolved \\
  handoff' = handoff \\
  tainted' = tainted
\end{schema}

\texttt{DeclineNote} is the person saying no to a held note, or the
confirmation gate dropping it. Only a held note can be declined, so a note
still in flight, one the gate has not yet seen, is never discarded this way.
The declined note leaves \texttt{pending} and is never recorded, so a
replacement the person confirms is a new note and the declined one cannot
make the handoff partial. A note still held when the call ends is not a
decline: it stays pending and the handoff is partial (H6).

\begin{schema}{SealHandoff}
  \Delta HandoffRecord \\
  \Xi Call
\where
  call \in \{ ended, failed \} \\
  handoff = collecting \\
  asked' = asked \\
  resolved' = resolved \\
  noted' = noted \\
  pending' = pending \\
  held' = held \\
  tainted' = tainted \\
  (call = ended \land endedBy = byAgent \land asked \subseteq resolved
    \land pending = \emptyset) \implies handoff' = complete \\
  \lnot (call = ended \land endedBy = byAgent \land asked \subseteq resolved
    \land pending = \emptyset) \implies handoff' = partial
\end{schema}

\texttt{SealHandoff} runs once the call has reached its final state, which
is after the last reply has played out, so that \texttt{call} is
\texttt{ended} or \texttt{failed} for good, and before
\texttt{SettleCall}. In the implementation that is the launcher's task
body after \texttt{call.run} returns or raises, not inside
\texttt{call.run}: after a remote hang-up the play-out can still turn the
call failed, and a seal taken before it could be \texttt{complete} for a
call that then fails. A call cancelled before its task ran never dialed;
the task's done callback seals it partial with no notes. The schema frames
\texttt{noted}, \texttt{resolved} and \texttt{pending}, which is H6.

The call operations above act on \texttt{Call}. In a briefed daemon each
one is lifted by \texttt{HandoffRecord} unchanged, except the ones that
start, receive or record a note, or seal.

\begin{zed}
  StartBriefed \defs (Start \land BeginHandoff) \land \Delta Briefed \\
  StartPlain \defs (Start \land BeginNoHandoff) \land \Delta Briefed \\
  ReceiveNoteB \defs (ReceiveNote \land \Xi Call) \land \Delta Briefed \\
  NoteDurable \defs (RunDurable \land NoteHandoff) \land \Delta Briefed \\
  HoldNoteB \defs (HoldNote \land \Xi Call) \land \Delta Briefed \\
  DeclineNoteB \defs (DeclineNote \land \Xi Call) \land \Delta Briefed \\
  ReplayNote \defs Replay \land \Xi HandoffRecord \\
  SealBriefed \defs SealHandoff \land \Delta Briefed \\
  ProvisionAgentB \defs ProvisionAgent \land \Xi HandoffRecord \land \Delta Briefed \\
  ConnectB \defs Connect \land \Xi HandoffRecord \land \Delta Briefed \\
  AgentEndCallB \defs AgentEndCall \land \Xi HandoffRecord \land \Delta Briefed \\
  EndB \defs End \land \Xi HandoffRecord \land \Delta Briefed \\
  SpeakB \defs Speak \land \Xi HandoffRecord \land \Delta Briefed \\
  InterruptB \defs Interrupt \land \Xi HandoffRecord \land \Delta Briefed \\
  FinishB \defs Finish \land \Xi HandoffRecord \land \Delta Briefed \\
  SettleCallB \defs SettleCall \land \Xi HandoffRecord \land \Delta Briefed
\end{zed}

\texttt{SettleCallB} has an invariant-derived precondition: it applies
only to a call whose handoff is not \emph{collecting} (H3), so the
implementation must seal first, after the final state. \texttt{StartBriefed} and
\texttt{StartPlain} replace the previous call's sealed handoff; the
previous handoff has been delivered in the status reply and the durable
file by then. A named stop in the briefed state is
\texttt{EndNamed}, then \texttt{SealBriefed}, then \texttt{SettleCallB}.

\section{Music suspension}

The call task holds music from before dialing through audio cleanup. Finishing,
failing, or cancelling the task releases only its own hold. A transport reset
(stop or source switch) clears the user's pause, not the call's hold.

\begin{schema}{MusicPause}
  userPaused, callHeld, paused : ZBOOL
\where
  paused = ztrue \iff (userPaused = ztrue \lor callHeld = ztrue)
\end{schema}

\begin{schema}{HoldMusic}
  \Delta MusicPause
\where
  callHeld = zfalse \land callHeld' = ztrue \\
  userPaused' = userPaused
\end{schema}

\begin{schema}{ReleaseMusic}
  \Delta MusicPause
\where
  callHeld = ztrue \land callHeld' = zfalse \\
  userPaused' = userPaused
\end{schema}

\begin{schema}{SetUserPause}
  \Delta MusicPause \\
  userPause? : ZBOOL
\where
  userPaused' = userPause? \\
  callHeld' = callHeld
\end{schema}

The pause-state effects of stopping and switching source are the same.

\begin{schema}{ResetMusicTransport}
  \Delta MusicPause
\where
  userPaused' = zfalse \\
  callHeld' = callHeld
\end{schema}

\begin{zed}
  MusicStop \defs ResetMusicTransport \\
  MusicSwitchSource \defs ResetMusicTransport
\end{zed}

\end{document}
